\DocumentMetadata{
	testphase={phase-III, math},
	tagging=on,
	tagging-setup={math/setup=mathml-SE},
	pdfstandard=ua-2,
	pdfversion=2.0,
	lang=en
}

\documentclass[11 pt]{article}
\usepackage[a4paper, top=15mm, bottom=15mm, left=20mm, right=20mm, footskip=8mm]{geometry}
\usepackage{hyperref}
\usepackage{float}
\usepackage{tikz}
\usetikzlibrary{calc}

\title{Another Document To Cross Reference}
\author{Henry Ginn}

\begin{document}

\def\title{Another Document To Cross-Reference}
\def\author{Henry Ginn}

\hypersetup{
	pdftitle={\title},
	pdfauthor={\author},
	pdfkeywords={Example},
	pdfdisplaydoctitle=true
}

\maketitle

This is an example document used to demonstrate how it is possible to cross-reference across documents using a combination of \LaTeX\ and Git submodule.

\section{Example Section The First}
\label{sec:ExampleSectionTheFirst}

This is a near-empty section included only to show that the following subsection number is correct when it is cross-referenced.

\section{Example Section The Second}
\label{sec:ExampleSectionTheSecond}

This is yet another near-empty section, this time included to show that the subsection number updates correctly.

\section{Example Section The Third}
\label{sec:ExampleSectionTheThird}

\subsection{Example Subsection}
\label{subsec:ExampleSubsection}

\begin{figure}[H]
	\centering
	\begin{tikzpicture}[scale=2.5, alt={A distorted grid with labelled cells and with some cells highlighted.}]
		\coordinate (00) at (1, 0);
\coordinate (40) at (4, 1);
\coordinate (44) at (5, 3.5);
\coordinate (04) at (0, 4);

\coordinate (10) at ($(00)!0.25!(40)$);
\coordinate (20) at ($(00)!0.50!(40)$);
\coordinate (30) at ($(00)!0.75!(40)$);
\coordinate (14) at ($(04)!0.25!(44)$);
\coordinate (24) at ($(04)!0.50!(44)$);
\coordinate (34) at ($(04)!0.75!(44)$);
\coordinate (01) at ($(00)!0.25!(04)$);
\coordinate (02) at ($(00)!0.50!(04)$);
\coordinate (03) at ($(00)!0.75!(04)$);
\coordinate (41) at ($(40)!0.25!(44)$);
\coordinate (42) at ($(40)!0.50!(44)$);
\coordinate (43) at ($(40)!0.75!(44)$);

\coordinate (11) at ($(10)!0.25!(14)$);
\coordinate (12) at ($(10)!0.50!(14)$);
\coordinate (13) at ($(10)!0.75!(14)$);
\coordinate (21) at ($(20)!0.25!(24)$);
\coordinate (22) at ($(20)!0.50!(24)$);
\coordinate (23) at ($(20)!0.75!(24)$);
\coordinate (31) at ($(30)!0.25!(34)$);
\coordinate (32) at ($(30)!0.50!(34)$);
\coordinate (33) at ($(30)!0.75!(34)$);

\coordinate (A00) at (02);
\coordinate (A02) at (04);
\coordinate (A100) at (42);
\coordinate (A102) at (44);

\coordinate (A10) at ($(A00)!0.1!(A100)$);
\coordinate (A20) at ($(A00)!0.2!(A100)$);
\coordinate (A30) at ($(A00)!0.3!(A100)$);
\coordinate (A40) at ($(A00)!0.4!(A100)$);
\coordinate (A50) at ($(A00)!0.5!(A100)$);
\coordinate (A60) at ($(A00)!0.6!(A100)$);
\coordinate (A70) at ($(A00)!0.7!(A100)$);
\coordinate (A80) at ($(A00)!0.8!(A100)$);
\coordinate (A90) at ($(A00)!0.9!(A100)$);
\coordinate (A12) at ($(A02)!0.1!(A102)$);
\coordinate (A22) at ($(A02)!0.2!(A102)$);
\coordinate (A32) at ($(A02)!0.3!(A102)$);
\coordinate (A42) at ($(A02)!0.4!(A102)$);
\coordinate (A52) at ($(A02)!0.5!(A102)$);
\coordinate (A62) at ($(A02)!0.6!(A102)$);
\coordinate (A72) at ($(A02)!0.7!(A102)$);
\coordinate (A82) at ($(A02)!0.8!(A102)$);
\coordinate (A92) at ($(A02)!0.9!(A102)$);

\coordinate (A01) at ($(A00)!0.5!(A02)$);
\coordinate (A41) at ($(A40)!0.5!(A42)$);
\coordinate (A61) at ($(A60)!0.5!(A62)$);
\coordinate (A101) at ($(A100)!0.5!(A102)$);

\coordinate (A11) at ($(A01)!0.1!(A101)$);
\coordinate (A21) at ($(A01)!0.2!(A101)$);
\coordinate (A31) at ($(A01)!0.3!(A101)$);
\coordinate (A51) at ($(A01)!0.5!(A101)$);
\coordinate (A71) at ($(A01)!0.7!(A101)$);
\coordinate (A81) at ($(A01)!0.8!(A101)$);
\coordinate (A91) at ($(A01)!0.9!(A101)$);

\fill[red!20] (A00) -- (A01) -- (A11) -- (A10);
\fill[red!20] (A91) -- (A101) -- (A102) -- (A92);
\fill[blue!20] (A01) -- (A02) -- (A12) -- (A11);
\fill[blue!20] (A90) -- (A100) -- (A101) -- (A91);

\draw (A00) -- (A40);
\draw (A01) -- (A41);
\draw (A02) -- (A42);
\draw (A60) -- (A100);
\draw (A61) -- (A101);
\draw (A62) -- (A102);
\draw (A00) -- (A02);
\draw (A10) -- (A12);
\draw (A20) -- (A22);
\draw (A30) -- (A32);
\draw (A70) -- (A72);
\draw (A80) -- (A82);
\draw (A90) -- (A92);
\draw (A100) -- (A102);

\node at (barycentric cs:A00=1,A01=1,A10=1,A11=1) {$L_1$};
\node at (barycentric cs:A10=1,A11=1,A20=1,A21=1) {$L_2$};
\node at (barycentric cs:A20=1,A21=1,A30=1,A31=1) {$L_3$};
\node at (barycentric cs:A01=1,A02=1,A11=1,A12=1) {$U_1$};
\node at (barycentric cs:A11=1,A12=1,A21=1,A22=1) {$U_2$};
\node at (barycentric cs:A21=1,A22=1,A31=1,A32=1) {$U_3$};
\node at (barycentric cs:A70=1,A71=1,A80=1,A81=1) {$L_{n-2}$};
\node at (barycentric cs:A80=1,A81=1,A90=1,A91=1) {$L_{n-1}$};
\node at (barycentric cs:A90=1,A91=1,A100=1,A101=1) {$L_n$};
\node at (barycentric cs:A71=1,A72=1,A81=1,A82=1) {$U_{n-2}$};
\node at (barycentric cs:A81=1,A82=1,A91=1,A92=1) {$U_{n-1}$};
\node at (barycentric cs:A91=1,A92=1,A101=1,A102=1) {$U_n$};
\node at (A51) {$\dots$};

	\end{tikzpicture}
	\caption[Example figure]{This is an example figure that I borrowed from my log of puzzles.}
	\label{fig:ExampleFigure}
\end{figure}

\end{document}